\section{PatchCore Enhancements}
\label{sec:enh}

We test three enhancements to PatchCore, each changing one factor relative to the WRN50 baseline under the same protocol.
All variants were specified \emph{a priori}, except one DINOv2 sensitivity run that was chosen after seeing test results and is marked post-hoc (\posthoc).
Table~\ref{tab:bonus} reports six-category means, and Fig.~\ref{fig:tradeoff} plots detection and localisation against latency.

\begin{table*}[t]
  \centering
  \caption{PatchCore enhancements, six-category means. Bold marks the best value among a-priori variants. \posthoc{}Post-hoc and test-informed; shown as a sensitivity check and not selectable. The input-pyramid and post-processing rows score images by the smoothed-map maximum and should be compared with the max-score control rather than with the Eq.~\eqref{eq:reweight} baseline.}
  \label{tab:bonus}
  \footnotesize
  \begin{tabular}{@{}ll@{\hspace{5pt}}r@{\hspace{5pt}}r@{\hspace{5pt}}r@{\hspace{5pt}}r@{\hspace{5pt}}r@{\hspace{5pt}}r@{}}
\toprule
Variant & Kind & \shortstack{Image\\AUROC} & \shortstack{Image F1\\(val thr)} & \shortstack{Pixel\\AUROC} & AU-PRO & \shortstack{Infer\\ms/img} & \shortstack{Peak\\GiB} \\
\midrule
WRN50 baseline (Eq.~7) & -- & 0.987 & 0.919 & 0.985 & 0.931 & 2.78 & 1.50 \\
\midrule
DINOv2 ViT-B/14 & a priori & 0.975 & 0.914 & 0.986 & 0.931 & 12.91 & 3.33 \\
DINOv2 ViT-S/14 & a priori & 0.968 & 0.928 & 0.986 & 0.929 & 6.33 & 2.22 \\
WRN50 + DINOv2-B & a priori & 0.965 & 0.906 & 0.983 & 0.914 & 13.77 & 3.08 \\
DINOv2-B no-LN\posthoc & post-hoc & 0.986 & 0.905 & 0.977 & 0.908 & 12.88 & 3.33 \\
\midrule
Layers 1+2+3 & a priori & 0.987 & 0.898 & \textbf{0.987} & \textbf{0.945} & 25.87 & 5.55 \\
Layers 2+3+4 & a priori & 0.922 & 0.812 & 0.933 & 0.706 & 3.25 & 1.67 \\
Max-score control & a priori & 0.988 & 0.934 & 0.985 & 0.931 & 2.79 & 1.50 \\
Pyramid mean & a priori & \textbf{0.990} & \textbf{0.943} & 0.987 & 0.940 & 12.81 & 3.45 \\
$z$-pyramid mean & a priori & \textbf{0.990} & \textbf{0.943} & 0.987 & 0.940 & 12.82 & 3.45 \\
\midrule
Top-1\% mean & a priori & 0.988 & 0.935 & 0.985 & 0.931 & 2.79 & 1.50 \\
Smoothed max & a priori & 0.988 & 0.934 & 0.985 & 0.931 & 2.79 & 1.50 \\
Smoothing $\sigma=0$ & a priori & 0.987 & 0.919 & 0.985 & 0.932 & 2.75 & 1.50 \\
Smoothing $\sigma=2$ & a priori & 0.987 & 0.919 & 0.985 & 0.932 & 2.78 & 1.50 \\
Smoothing $\sigma=8$ & a priori & 0.987 & 0.919 & 0.985 & 0.923 & 2.81 & 1.50 \\
Global threshold $z=3$ & a priori & 0.987 & 0.890 & 0.985 & 0.931 & 2.79 & 1.50 \\
Global pooled val & a priori & 0.987 & 0.869 & 0.985 & 0.931 & 2.79 & 1.50 \\
\bottomrule
\end{tabular}

\end{table*}

\textbf{DINOv2 backbone.}
We replace WRN50 by self-supervised DINOv2 ViT-B/14 or ViT-S/14 features~\cite{oquab2024dinov2}, or concatenate the two backbones, motivated by the strong few-shot results of AnomalyDINO~\cite{damm2025anomalydino}.
Textures improve: mean texture F1 rises from 0.931 to 0.987--0.993 and texture AU-PRO from 0.914 to 0.941--0.945.
However, all three a-priori variants \emph{lower} mean image AUROC (\bnDinoBImgAuroc{}, \bnDinoSImgAuroc{}, \bnWrnDinoImgAuroc{} vs \bnBaseImgAuroc{}), driven by screw (0.971$\to$0.809--0.879), at 2.3--4.9$\times$ the inference time.
A post-hoc diagnostic on screw traced the regression mainly to processing rather than to the features themselves.
Applying DINOv2's \emph{final} LayerNorm to intermediate blocks concentrates about 40\% of the token variance in 5 of 768 channels, which then dominate Euclidean nearest-neighbour distances.
The CNN-oriented 3$\times$3 patch aggregation also hurts.
Removing the LayerNorm recovers mean image AUROC to \bnDinoNoLNImgAuroc{} but lowers texture AU-PRO.
Because this variant was chosen after inspecting test results, we report it only as a sensitivity check.

\textbf{Multi-scale fusion.}
Adding \texttt{layer1} to the feature hierarchy (L1+2+3) gives the best localisation: AU-PRO rises from \bnBaseAupro{} to \bnMsLOneTwoThreeAupro{}, with gains in every category (e.g.\ +0.025 on tile).
It costs image F1 (\bnMsLOneTwoThreeImgFone{} vs \bnBaseImgFone{}), $9.3\times$ inference time and 5.55\,GiB of memory.
Finer stride-4 features help the sub-cell and boundary failures of Section~\ref{sec:failures}, but they also add rare-but-normal peaks that disturb the image decision.
Adding \texttt{layer4} instead hurts badly (AU-PRO \bnMsLTwoThreeFourAupro{}), because coarse 7$\times$7 semantics dilute fine defect evidence.
An \emph{input pyramid} fuses independently fitted banks at scales 1.0 and 1.5 by averaging their maps.
Against the max-score control, it improves image AUROC (\bnPyrMeanImgAuroc{} vs \bnMsMaxCtrlImgAuroc{}), image F1 (\bnPyrMeanImgFone{} vs \bnMsMaxCtrlImgFone{}) and AU-PRO (\bnPyrMeanAupro{} vs \bnMsMaxCtrlAupro{}), at $4.6\times$ the inference time.
It is the best a-priori variant for detection.
Normalising each scale by training-only z-scores before averaging makes no difference.

\textbf{Post-processing.}
Replacing the Eq.~\eqref{eq:reweight} image score with the maximum or the top-1\% mean of the smoothed map leaves AUROC almost unchanged but raises mean F1 from \bnBaseImgFone{} to \bnPpMaxImgFone{} and \bnPpTopkImgFone{}.
The aggregation therefore changes calibration, not ranking.
Its effect is category-dependent: top-1\% helps capsule and tile but hurts screw.
Smoothing with $\sigma\le4$ makes no practical difference (AU-PRO within 0.0016), whereas $\sigma=8$ lowers AU-PRO in every category.
Finally, replacing per-category thresholds with \emph{one global threshold}, either $z=3$ on validation-normalised scores or a 99\% quantile of the pooled validation scores, lowers mean F1 to \bnPpGlobalZImgFone{} and \bnPpGlobalPooledImgFone{}.
Per-product calibration therefore remains necessary.

\begin{figure}[t]
  \centering
  \includegraphics[width=\columnwidth]{fig_tradeoff.pdf}
  \caption{Accuracy--latency trade-off of the PatchCore model variants (six-category means; inference time on a log$_2$ axis): (a)~image AUROC, (b)~AU-PRO. Dotted lines mark the WRN50 baseline; the max-score control nearly coincides with it ($|\Delta|\leq0.001$). All values are single-seed estimates. The hollow marker (\dag) is the post-hoc DINOv2 run. Post-processing variants are omitted because they do not change latency (Table~\ref{tab:bonus}). A value below the axis range is drawn on the axis edge with its value.}
  \label{fig:tradeoff}
\end{figure}
